\begin{definition}[title = The Church-Turing Thesis]
    \begin{itemize}
        \renewcommand{\labelitemi}{-}
        \item Every physically realizable computer can be simulated by a Turing Machine with a overhead in time.
        \item The class of computable tasks would not be larger if formalized in a realistic way, but differently.
    \end{itemize}
\end{definition}
\begin{definition}[title = The Strong Church-Turing Thesis]
    \begin{itemize}
        \renewcommand{\labelitemi}{-}
        \item Every physically realizable computer can be simulated by a Turing Machine with \textbf{polynomial} overhead in time ($n$ steps on the computer requires $n^c$ on Turing Machines, where $c$ only depends on the computer), and viceversa.
        \item The class \textbf{P} would be \textbf{the same} if defined based on other realistic models of computation.
    \end{itemize}
\end{definition}

Note: by the second thesis we specify that for any task belonging to the complexity class \textbf{P} its resolution by any Turing Machine always requires
a \textit{polynomial} overhead in time.